% !TeX root = ../main.tex
\chapter{How Check2HGI and the Joint Model Work}
\label{apx:check2hgi-joint-model}

This appendix explains the two methods used in the final study as an end-to-end
system. The emphasis is practical: what information enters each stage, how it is
transformed, what is learned, and what must be reproduced to obtain the same
experimental setting. Equations are included only where prose alone would leave
an implementation choice ambiguous.

The system has two separately trained parts. Check2HGI first learns one vector for
each check-in and one vector for each geographic region without using the labels of
the two forecasting tasks. The joint model then consumes sequences of those vectors
and learns to predict both the category and the region of the next visit. The first
part answers ``how should a visit and its geographic context be represented?'' The
second answers ``how can the two forecasting tasks exchange useful information
without forcing them to use exactly the same predictor?''

\section{The complete pipeline at a glance}
\label{sec:apx-methods-pipeline}

Table~\ref{tab:apx-methods-flow} summarizes the complete data flow. The two training
stages are intentionally separated. Check2HGI is trained once to produce reusable
representations. Those representations are then treated as fixed inputs to the
supervised joint model.

\begin{longtable}{p{0.07\textwidth}p{0.23\textwidth}p{0.37\textwidth}p{0.22\textwidth}}
\caption{End-to-end flow from raw check-ins to the two next-visit predictions.}
\label{tab:apx-methods-flow}\\
\toprule
\textbf{Stage} & \textbf{Input} & \textbf{Operation} & \textbf{Output} \\
\midrule
\endfirsthead
\toprule
\textbf{Stage} & \textbf{Input} & \textbf{Operation} & \textbf{Output} \\
\midrule
\endhead
1 & Raw check-in records and region polygons & Clean identifiers, sort each user's visits by time, assign places to polygons, and create stable integer mappings & Validated check-ins, places, regions, users, and categories \\
2 & Validated records & Build the check-in succession graph, the Delaunay place graph, the polygon-adjacency graph, and the check-in--place--region--city hierarchy & A four-level heterogeneous mobility graph \\
3 & Graph and node features & Train Check2HGI with three contrastive hierarchy losses and two label-free auxiliary losses & Learned check-in, place, and region vectors \\
4 & Ordered user histories and learned vectors & Create sliding windows containing nine observed visits and use visit ten as the target & Paired category and region training examples \\
5 & Two $9\times64$ input sequences & Encode each stream, exchange context through two cross-attention blocks, and preserve a private spatial route for the region task & Shared-context and private task features \\
6 & Task features and next-visit labels & Optimize one weighted objective and select checkpoints using both tasks & Seven category logits and one logit per region class \\
\bottomrule
\end{longtable}

The important separation is between \emph{representation learning} in Stages 2--3
and \emph{forecast learning} in Stages 4--6. Check2HGI does not receive the next
category or next region as a training target. Conversely, the joint model does not
rebuild the graph. It reads the vectors exported by Check2HGI.

\section{From a check-in table to the Check2HGI graph}
\label{sec:apx-check2hgi-input}

\subsection{Required record and stable ordering}

Each observation is written as
\begin{equation}
    c_i=(u_i,p_i,l_i,g_i,t_i).
    \label{eq:apx-checkin-tuple}
\end{equation}
Here, $c_i$ is check-in $i$, $u_i$ is its user, $p_i$ is its place, $l_i$
contains latitude and longitude, $g_i$ is one of the seven place categories, and
$t_i$ is the timestamp. A spatial join between $l_i$ and the polygon file derives
the region. Places outside every polygon are removed. Identifiers are remapped to
contiguous integers so that each row of an embedding matrix has one unambiguous
entity mapping.

Visits are sorted by user and timestamp before any temporal edge or forecasting
window is created. This ordering is part of the method, not a presentation detail.
Changing it changes both the graph neighborhoods seen by Check2HGI and the examples
seen by the joint model.

\subsection{Four node levels and three kinds of structure}

The hierarchy contains check-in, place, region, and city nodes. Formally, it is a
heterogeneous graph $\mathcal{G}=(\mathcal{V},\mathcal{E},\mathcal{T}_v,
\mathcal{T}_e)$ with node-type mapping $\phi:\mathcal{V}\rightarrow\mathcal{T}_v$
and edge-type mapping $\psi:\mathcal{E}\rightarrow\mathcal{T}_e$. The set
$\mathcal{V}$ contains every node, $\mathcal{E}$ contains every edge,
$\mathcal{T}_v$ contains the four node types, and $\mathcal{T}_e$ identifies the
relations described below.

The graph combines three different notions of relatedness.

\begin{enumerate}
    \item \textbf{Temporal succession at the check-in level.} Consecutive visits by
    the same user are connected in both directions. A short gap receives a larger
    weight than a long gap. Before global min--max rescaling, the weight is
    \begin{equation}
        \widetilde{a}_{ij}=\exp\!\left(-\frac{|t_j-t_i|}{3600}\right).
        \label{eq:apx-temporal-weight}
    \end{equation}
    Here, $\widetilde{a}_{ij}$ is the raw weight between consecutive check-ins
    $c_i$ and $c_j$, $|t_j-t_i|$ is their time difference in seconds, and $3600$
    is the one-hour decay constant. If these raw weights are not all equal, they
    are rescaled to the interval $[0,1]$ over the complete check-in graph.

    \item \textbf{Spatial proximity at the place and region levels.} Places are
    joined by a Delaunay triangulation of their coordinates. Regions are joined
    when their polygons intersect. These two graphs allow local spatial context to
    move between nearby places and neighboring regions.

    \item \textbf{Geographic membership across levels.} Every check-in belongs to
    one place, every retained place belongs to one region, and every region belongs
    to the city. The implementation stores these memberships as assignment vectors
    and uses segmented pooling rather than constructing a large sparse matrix for
    every cross-level relation.
\end{enumerate}

Category co-occurrence is not an edge relation in the final configuration.
Category enters through node features and an auxiliary reconstruction target.
Likewise, the active check-in graph contains user-succession edges only. Repeated
visits to the same place meet through their common place node; they are not joined
by an additional same-place check-in edge.

\subsection{What a check-in node initially knows}

Each check-in starts with 11 values: a seven-dimensional one-hot category vector,
two cyclic values for hour of day, and two cyclic values for day of week. Sine and
cosine encodings make midnight adjacent to 23:00 and Sunday adjacent to Monday in
feature space. Raw latitude and longitude are not appended to this feature vector.
Coordinates instead determine polygon membership and the spatial graphs. This
distinction prevents the representation from treating geographic coordinates as
ordinary categorical or temporal features.

\section{How Check2HGI transforms the graph}
\label{sec:apx-check2hgi-forward}

Check2HGI performs a bottom-up pass through the hierarchy. Every level has a
specific role: graph convolutions mix information between neighbors at the same
level, while attention pooling summarizes a variable number of lower-level nodes
into one upper-level vector.

\subsection{Step 1: contextualize each check-in}

Two weighted graph-convolution layers operate on the temporal succession graph.
The first maps the 11 input features to 64 dimensions, followed by layer
normalization, PReLU, and dropout $0.1$. The second maps $64$ to $64$ and is added
back to the first-layer state as a residual connection. The result for check-in
$i$ is $\mathbf{x}_i\in\mathbb{R}^{64}$.

This operation changes the meaning of a check-in vector. It no longer describes
only the visit's own category and time. It also reflects the visits immediately
before and after it in the same user's trajectory, with the temporal edge weights
controlling how strongly those neighbors contribute.

\subsection{Step 2: summarize visits at each place}

All check-in vectors assigned to a place are pooled with four attention heads. The
pooling is not a simple mean. It learns which visits are more informative for
describing that place. A visit during an unusual time period, for example, can
receive a different weight from a routine visit even though both refer to the same
coordinates.

The output is a semantic place embedding $\mathbf{e}_{p}\in\mathbb{R}^{64}$ for
each place $p$. Four heads divide the 64-dimensional representation into four
16-dimensional views. A learned seed query asks the same general question of every
place, while the keys and values derived from its visits determine the weighted
answer. A residual projection and PReLU produce the final place vector.

\subsection{Step 3: add the spatial place neighborhood}

The place-to-region route combines the semantic place embedding with a trainable
64-dimensional place table initialized from POI2Vec. A learned scalar, initialized
to $1.0$, controls the initialized table's contribution. One weighted
$64\rightarrow64$ graph convolution then propagates information over the Delaunay
place graph, followed by channelwise PReLU.

The semantic place embedding is detached on this spatial route. Consequently, the
higher-level place--region and region--city losses cannot rewrite the check-in
encoder through this branch. The lower level still learns through its own
check-in--place objective. This stop-gradient boundary gives the spatial hierarchy
a stable place basis while retaining end-to-end learning inside each intended
route.

\subsection{Step 4: summarize places into regions and regions into the city}

Places within a region are pooled by another four-head attention module. A
$64\rightarrow64$ graph convolution then exchanges information between adjacent
region polygons. The resulting row $\mathbf{z}_r\in\mathbb{R}^{64}$ is the learned
embedding of region $r$.

The city vector is the sigmoid of an area-weighted sum of all region vectors. The
polygon areas are not normalized to sum to one. This city vector is used while
training the highest hierarchy boundary; downstream forecasting consumes the
check-in and region vectors rather than the city vector.

\section{How Check2HGI learns without forecast labels}
\label{sec:apx-check2hgi-learning}

Check2HGI follows the graph-infomax principle of making representations from true
relationships score higher than representations from corrupted relationships
\cite{velickovic2019deep,huang2023hgi}. It applies this idea at three boundaries.

\begin{description}
    \item[Check-in to place.] A check-in vector is paired with its true place
    vector as a positive example and with a randomly selected different place as a
    negative example.

    \item[Place to region.] A place vector is paired with its true region as a
    positive example and a different region as a negative example. For graphs with
    fewer than $50{,}000$ places, 25\% of negative-sampling attempts search for a
    moderately similar region, with category-histogram cosine similarity strictly
    between $0.6$ and $0.8$. This produces a harder comparison. The similarity is
    used only for negative selection, not as a graph edge.

    \item[Region to city.] Real region vectors are compared with the city summary.
    Negative regions are produced by independently permuting the rows of the
    check-in feature matrix and running a second forward pass while keeping the
    temporal topology fixed.
\end{description}

Two auxiliary objectives regularize the hierarchy. The masked-place objective
zeros 15\% of the semantic place vectors, averages their Delaunay neighbors, and
uses a $64\rightarrow128\rightarrow7$ PReLU decoder to reconstruct the masked
place's empirical category distribution. Its error is the cube of a clipped cosine
distance. The place-table anchor penalizes movement away from the initial POI2Vec
table with mean squared error.

The implemented representation objective is
\begin{equation}
\begin{split}
    \mathcal{L}_{\mathrm{Check2HGI}}={}&
    0.4\mathcal{L}_{C\!P}+0.3\mathcal{L}_{P\!R}
    +0.3\mathcal{L}_{R\!\Omega}\\
    &+0.3\mathcal{L}_{\mathrm{mp}}+0.1\mathcal{L}_{\mathrm{anc}}.
\end{split}
\label{eq:apx-check2hgi-loss}
\end{equation}
Here, $\mathcal{L}_{\mathrm{Check2HGI}}$ is the scalar loss used to train the
representation model; $\mathcal{L}_{C\!P}$, $\mathcal{L}_{P\!R}$, and
$\mathcal{L}_{R\!\Omega}$ are the contrastive losses at the check-in--place,
place--region, and region--city boundaries; $\mathcal{L}_{\mathrm{mp}}$ is the
masked-place reconstruction loss; and $\mathcal{L}_{\mathrm{anc}}$ is the
place-table anchor. The coefficients are additive design weights and are not
required to sum to one.

Training uses full-batch Adam for 500 epochs with learning rate $10^{-3}$, no
weight decay, default Adam coefficients, and global gradient-norm clipping at
$0.9$. There is no learning-rate decay or early stopping. The saved state is the
epoch with the smallest complete training loss. Seed 42 is used for this stage.

\section{What Check2HGI exports}
\label{sec:apx-check2hgi-output}

The representation stage exports separate lookup tables for check-ins, places, and
regions, together with the identifier maps needed to retrieve the correct row. The
joint model uses two of those outputs.

\begin{enumerate}
    \item The category stream looks up the 64-dimensional check-in vector
    $\mathbf{x}_i$ produced by the temporal graph encoder.
    \item The region stream looks up the 64-dimensional region vector
    $\mathbf{z}_{r_i}$ associated with the visited place's polygon.
\end{enumerate}

The two vectors remain separate. In particular, the method does not concatenate a
place embedding, raw coordinates, and temporal features into one 128-dimensional
visit vector. Category and cyclic time are fused by the temporal graph encoder;
coordinates shape the spatial hierarchy; and the learned hierarchy losses connect
the levels during representation training.

For each user $u$, a forecasting example contains the ordered check-in matrix
\begin{equation}
    \mathbf{H}_u=(\mathbf{x}_1,\mathbf{x}_2,\ldots,\mathbf{x}_k)
    \in\mathbb{R}^{k\times d'},\qquad k=9,\quad d'=64.
    \label{eq:apx-user-sequence}
\end{equation}
Here, $\mathbf{H}_u$ is the category-stream history, $\mathbf{x}_j$ is the
Check2HGI vector for historical visit $j$, $k$ is the number of observed visits,
and $d'$ is the check-in-vector width. A second matrix of the same shape contains
the corresponding region vectors. All-zero vectors denote padded positions, and
the model converts those positions into attention masks.

Windows have stride one. Visits 1--9 predict visit 10, visits 2--10 predict visit
11, and so forth. Users need at least ten visits. No synthetic final target is
added. User identities, not individual windows, determine the five folds, so one
person's windows cannot appear in both the training and evaluation partitions of a
fold.

The representation model is trained on the complete graph before the supervised
folds are fitted. The evaluation is therefore user-disjoint for the forecasting
labels, but the label-free representation stage is transductive with respect to
the graph. Reproducing the protocol requires preserving this distinction.

\section{How the joint model processes the two histories}
\label{sec:apx-joint-forward}

The joint architecture is not one shared encoder followed by two small output
layers. It combines private encoders, an interaction subsystem, and task-specific
heads. Information is shared through cross-attention, while each task retains a
route suited to its own output.

\subsection{Step 1: encode the two modalities separately}

The check-in and region histories enter independent encoders with the same shape
but different parameters. Each encoder applies three linear transformations:
$64\rightarrow256$, $256\rightarrow256$, and $256\rightarrow256$. Every
transformation is followed by ReLU and layer normalization. Dropout $0.1$ follows
the first two transformations but not the third. The output consists of two
$9\times256$ sequences.

Keeping these encoders separate matters because a check-in vector and a region
vector do not represent the same object. Equal tensor widths do not imply equal
semantics or shared weights.

\subsection{Step 2: exchange context with cross-attention}

Two bidirectional cross-attention blocks connect the streams. Within each block,
the category stream first queries the region stream. The region stream then queries
the already updated category stream. Each direction has its own projections,
normalizations, and feed-forward network.

For one attention head, the core operation is
\begin{equation}
    \operatorname{Attn}(\mathbf{Q},\mathbf{K},\mathbf{V})
    =\operatorname{softmax}\!\left(
      \frac{\mathbf{Q}\mathbf{K}^{\mathsf{T}}}{\sqrt{d_h}}
    \right)\mathbf{V}.
    \label{eq:apx-cross-attention}
\end{equation}
Here, $\mathbf{Q}$ contains query vectors from the stream being updated,
$\mathbf{K}$ and $\mathbf{V}$ contain key and value vectors from the other stream,
$d_h=64$ is the head width, and $\operatorname{softmax}$ converts compatibility
scores into weights over valid historical positions. Padded keys are masked before
the softmax.

Each cross-attention module uses four heads, model width 256, and attention dropout
$0.15$. A residual connection and layer normalization follow each attention
operation. Each stream then passes through its own
$256\rightarrow256\rightarrow256$ feed-forward network with GELU and dropout
$0.15$, followed by another residual connection and normalization. After two
blocks, separate final layer normalizations produce the category-context and
region-context sequences.

The interaction stack is optimized as the shared subsystem, but its primitive
weights are directional rather than tied. Sharing occurs because each task reads
the other task's activations and both losses update the coupled system. This is
more precise than describing the model as conventional hard parameter sharing.

\subsection{Step 3: predict the next category}

The category-context sequence enters a four-layer unidirectional gated recurrent
unit (GRU). Its input and hidden widths are 256, and dropout $0.1$ is applied
between GRU layers. The model selects the top-layer hidden state at the final
nonpadding position, applies layer normalization and dropout $0.1$, and maps the
result to seven logits. The largest logit identifies the predicted next category.

The active implementation uses four GRU layers. This depth is set when the joint
model instantiates the category head and is therefore the value required for
reproduction.

\subsection{Step 4: predict the next region with two routes}

The region head deliberately preserves both a private spatial route and a route
that has interacted with category.

\begin{description}
    \item[Private route.] The raw $9\times64$ region sequence is projected to width
    128 and processed by a STAN-style spatio-temporal attention tower
    \cite{luo2021stan}. It uses four attention heads, dropout $0.3$, learned
    $4\times9\times9$ pairwise-position biases, and a
    $128\rightarrow512\rightarrow128$ GELU feed-forward network.

    \item[Shared-context route.] The $9\times256$ region-context sequence from the
    cross-attention stack is projected to width 128 and processed by a separate
    tower. This tower uses eight attention heads, dropout $0.1$, learned
    $8\times9\times9$ pairwise-position biases, and the same feed-forward widths.
\end{description}

Each tower uses one self-attention stage to contextualize the history and a second
attention stage to retrieve the final valid query. Both therefore produce a
128-dimensional vector. The private vector is added to a linear projection of the
shared-context vector, multiplied by a trainable scalar initialized to $0.1$.
Layer normalization, dropout $0.1$, and a final linear layer then produce one logit
for each region in the dataset-specific region map.

The code contains an optional region-transition prior, but its coefficient is fixed
at zero in the final model. It contributes nothing to the logits. The reported
model also uses no knowledge-distillation term.

\section{How the two tasks are trained together}
\label{sec:apx-joint-training}

Both task heads use ordinary mean cross-entropy. Their exact combination is
\begin{equation}
    \mathcal{L}_{\mathrm{total}}
    =0.75\mathcal{L}_{C}+0.25\mathcal{L}_{R}.
    \label{eq:apx-joint-loss}
\end{equation}
Here, $\mathcal{L}_{\mathrm{total}}$ is the scalar loss used for one backward pass,
$\mathcal{L}_{C}$ is mean cross-entropy for the seven next-category classes, and
$\mathcal{L}_{R}$ is mean cross-entropy over the dataset's region classes. The
weights are fixed. No uncertainty weighting, GradNorm, gradient surgery, class
weights, label smoothing, or dynamic task balancing is active.

AdamW uses three parameter groups. The category group contains the check-in encoder
and GRU head. The region group contains the region encoder and both region towers.
The shared group contains the two cross-attention blocks and final stream
normalizations. All groups use weight decay $0.05$, Adam coefficients $(0.9,0.999)$,
and numerical constant $10^{-8}$. The category peak learning rate is $10^{-3}$ and
the region peak is $3\times10^{-3}$. The shared peak is $3\times10^{-3}$ for
Alabama, Arizona, and Florida, and $10^{-3}$ for California, Texas, and Istanbul.
A per-group OneCycle schedule uses the default warmup fraction $0.3$ and advances
once after every optimizer step.

The category and region data loaders use the same user-disjoint fold but shuffle
independently during training. The shorter loader is cycled until the longer loader
is exhausted. A category batch and a region batch in the same optimizer step can
therefore contain different windows. Validation uses aligned examples. This pairing
rule is easy to overlook and must be retained for an exact reproduction.

Training uses batch size 8192 for each task loader, 50 epochs, no gradient
accumulation, and global gradient-norm clipping at $1.0$. Early stopping is disabled.
Experiments use seeds $\{0,1,7,100\}$ and five user-disjoint folds. For each run,
the selected checkpoint maximizes the geometric mean of category macro-F1 and
region Accuracy@10, so model selection cannot improve one task while ignoring the
other.

\section{Compact implementation specification}
\label{sec:apx-methods-specification}

Tables~\ref{tab:apx-check2-spec} and~\ref{tab:apx-joint-spec} collect the settings
that determine tensor compatibility and training behavior. The number of region
classes is deliberately absent as a fixed value because it is read from each
dataset's region map.

\begin{longtable}{p{0.30\textwidth}p{0.60\textwidth}}
\caption{Check2HGI implementation specification.}
\label{tab:apx-check2-spec}\\
\toprule
\textbf{Component} & \textbf{Implemented setting} \\
\midrule
\endfirsthead
\toprule
\textbf{Component} & \textbf{Implemented setting} \\
\midrule
\endhead
Check-in input & 7 category indicators plus 4 cyclic time features; width 11 \\
Check-in encoder & Two weighted graph convolutions, $11\rightarrow64\rightarrow64$; residual second layer \\
Check-in nonlinearity & Layer normalization, shared-slope PReLU, and dropout $0.1$ after the first layer \\
Check-in--place pooling & Four attention heads; 16 values per head; output width 64 \\
Spatial place route & POI2Vec-initialized trainable table; scalar initialized to $1.0$; one weighted $64\rightarrow64$ Delaunay graph convolution; channelwise PReLU \\
Place--region pooling & Four attention heads; output width 64 \\
Region graph & One unweighted $64\rightarrow64$ convolution over polygon adjacency; channelwise PReLU \\
Masked-place decoder & Mask probability $0.15$; neighbor mean; $64\rightarrow128\rightarrow7$ with PReLU \\
Loss weights & $0.4$ check-in--place, $0.3$ place--region, $0.3$ region--city, $0.3$ masked-place, $0.1$ table anchor \\
Optimization & Full-batch Adam; learning rate $10^{-3}$; weight decay 0; 500 epochs; clip norm $0.9$; seed 42 \\
\bottomrule
\end{longtable}

\begin{longtable}{p{0.30\textwidth}p{0.60\textwidth}}
\caption{Joint-model implementation specification.}
\label{tab:apx-joint-spec}\\
\toprule
\textbf{Component} & \textbf{Implemented setting} \\
\midrule
\endfirsthead
\toprule
\textbf{Component} & \textbf{Implemented setting} \\
\midrule
\endhead
Window construction & Nine observed visits, tenth visit as target, stride 1, minimum user history 10 \\
Input tensors & Separate check-in and region sequences, each $9\times64$; zero-vector padding \\
Private encoders & Independent $64\rightarrow256\rightarrow256\rightarrow256$ linear stacks with ReLU and layer normalization; dropout $0.1$ after the first two transformations \\
Interaction subsystem & Two bidirectional cross-attention blocks; width 256; four heads; attention and feed-forward dropout $0.15$ \\
Interaction feed-forward network & Independent per stream and block; $256\rightarrow256\rightarrow256$ with GELU \\
Category head & Four-layer unidirectional GRU; input and hidden width 256; interlayer dropout $0.1$; seven-class output \\
Private region tower & Raw region input; width 128; four heads; dropout $0.3$; feed-forward width 512 \\
Shared-context region tower & Cross-attention input; width 128; eight heads; dropout $0.1$; feed-forward width 512 \\
Region fusion & Private feature plus projected shared feature scaled by a trainable scalar initialized to $0.1$ \\
Active supervised loss & $0.75$ category cross-entropy plus $0.25$ region cross-entropy \\
Optimization & AdamW; weight decay $0.05$; OneCycle; batch size 8192 per task; 50 epochs; clip norm $1.0$ \\
Peak learning rates & Category $10^{-3}$; region $3\times10^{-3}$; shared $3\times10^{-3}$ for AL/AZ/FL and $10^{-3}$ for CA/TX/Istanbul \\
Evaluation design & Five user-disjoint folds; seeds $\{0,1,7,100\}$; selection by geometric mean of category macro-F1 and region Accuracy@10 \\
Inactive options & Transition-prior coefficient 0; knowledge-distillation weight 0; no class weighting or early stopping \\
\bottomrule
\end{longtable}

\section{Reproduction procedure and sanity checks}
\label{sec:apx-methods-reproduction}

An independent implementation can reproduce the method with the following order of
operations.

\begin{enumerate}
    \item Validate the required check-in fields, sort visits stably by user and
    timestamp, spatially join places to the region polygons, remove unmatched
    places, and freeze the identifier maps.

    \item Create 11-dimensional check-in features. Build bidirectional consecutive
    user edges with the one-hour exponential decay, a symmetrized Delaunay place
    graph, polygon-adjacency region edges, and the three hierarchy assignments.

    \item Initialize the place table from POI2Vec and train Check2HGI for 500
    full-batch epochs with Equation~\ref{eq:apx-check2hgi-loss}. Export the best-loss
    check-in, place, and region tables with their mappings.

    \item Rejoin exported rows to the chronologically ordered records. Check that
    every check-in and region lookup has width 64 and that each row resolves through
    exactly one stable identifier.

    \item Generate stride-one windows. For every eligible position, store nine
    historical check-in vectors, nine historical region vectors, the tenth visit's
    category, and the tenth visit's region. Do not append a fabricated tail example.

    \item Assign complete users to five folds and apply the same fold assignment to
    both tasks. Instantiate the joint model with seven category outputs and the
    region count read from that dataset's region map.

    \item Train for 50 epochs using Equation~\ref{eq:apx-joint-loss}, the three
    AdamW groups, the dataset-specific shared learning rate, independently shuffled
    task loaders, and maximum-size cycling. Save the checkpoint selected by the two
    task metrics together.

    \item Repeat the five folds for the required seeds. Report category macro-F1 and
    region Accuracy@10 from user-disjoint test partitions and retain the fold and
    seed identifiers in every result row.
\end{enumerate}

Before accepting a reproduction, the following checks should pass.

\begin{itemize}
    \item The active check-in graph has succession edges only; category
    co-occurrence and same-place check-in edges are absent.
    \item The representation dimensions are $11\rightarrow64$ at the check-in
    encoder and 64 at the place and region outputs.
    \item The joint model receives two separate $9\times64$ tensors, not one
    concatenated $9\times128$ tensor.
    \item The category GRU has four layers, and the cross-attention subsystem has
    two blocks with four heads per direction.
    \item The private region tower reads the raw region sequence. The shared-context
    tower reads the cross-attention output. Their 128-dimensional results are added,
    not concatenated.
    \item The transition-prior coefficient is zero, and the only supervised loss is
    the fixed $0.75/0.25$ combination.
    \item User identities are disjoint across supervised folds even though the
    label-free representation model was trained on the complete graph.
\end{itemize}

\section{Implementation map}
\label{sec:apx-methods-implementation-map}

The principal implementation points are listed below so that the conceptual stages
can be matched to executable definitions.

\begin{description}
    \item[Graph preparation and features:] \path{research/embeddings/check2hgi/preprocess.py}
    \item[Check-in encoder and masked-place decoder:] \path{research/embeddings/check2hgi/model/variants.py}
    \item[Check-in--place attention:] \path{research/embeddings/check2hgi/model/Checkin2POI.py}
    \item[Hierarchy, discriminators, and representation loss:] \path{research/embeddings/check2hgi/model/Check2HGIModule.py}
    \item[Region pooling and adjacency convolution:] \path{research/embeddings/hgi/model/RegionEncoder.py}
    \item[Task encoders and joint model assembly:] \path{src/models/mtl/mtlnet/model.py}
    \item[Bidirectional cross-attention:] \path{src/models/mtl/mtlnet_crossattn/model.py}
    \item[Dual-tower region routing:] \path{src/models/mtl/mtlnet_crossattn_dualtower/model.py}
    \item[Category GRU head:] \path{src/models/next/next_gru/head.py}
    \item[Region attention towers and fusion:] \path{src/models/next/next_stan_flow_dualtower/head.py}
    \item[Static multitask objective:] \path{src/losses/static_weight/loss.py}
    \item[Optimizer groups and scheduling:] \path{src/training/helpers.py}
\end{description}

Taken together, these stages give the whole methodological picture. Check2HGI turns
individual visits into representations that are simultaneously temporal,
place-aware, and region-aware. The joint model then lets category and region
histories inform one another while preserving the private spatial path needed for
fine-grained region ranking. The architecture's central idea is therefore not only
to share parameters. It is to control where information is fused, where it remains
private, and which training signal is allowed to update each route.
